% Standalone section fragment for inclusion in the same-action momentum note.
% Requires amsmath and amssymb. No independent preamble or bibliography.
\section{Exact tensor-action certification of the transverse projectors}
\label{sec:exact-transverse-action-certificate}

The earlier integer projectors were exact as projectors, but their
identification with the scalar action was only tested numerically.  We now
close that particular gap by an independent Gaussian-integer contraction
calculation.  No floating Hessian, eigensolver, tolerance, or rational fit to
a Hessian enters this certificate.  The claim remains a theorem about the
declared real scalar action on its fixed-orientation radial slice, not a
complete gauged pole or vacuum-stability theorem.

\subsection{A rational background and the canonical five-form map}

Let $q=(q_1<\cdots<q_5)$ label the $252$ five-element subsets of
$\{0,\ldots,9\}$.  Pair $q$ with its complement $q^c$ in the same order
used in the action.  If $s_q$ is the sign of $(q,q^c)$ and the chosen
Hodge eigenvalue is $d\,i$, $d\in\{1,-1\}$, the canonical real
$126+126$ coordinates obey
\begin{align}
 \Sigma_q&=\frac{x_k+i x_{k+126}}2,\\
 \Sigma_{q^c}&=-d\,i\,s_q\frac{x_k+i x_{k+126}}2.
\end{align}
Thus $C_j=2\,\partial\Sigma/\partial x_j$ has Gaussian-integer
quint coefficients.  No $\sqrt2$ remains in this canonical Sigma map.

Choose
\begin{equation}
 (w,\sigma,v_s)=(5,\sqrt2,\sqrt2),\qquad
 \Phi=\operatorname{diag}(-2,-2,-2,-2,-2,-2,3,3,3,3),\quad S=1.
\end{equation}
Writing $B=\Sigma(x_0)$, the tensor $A=4B$ has coefficient
$i^{n_o(q)}$ when $q$ selects one index from each ordered pair
$(0,1),(2,3),\ldots,(8,9)$, and zero otherwise; $n_o(q)$ counts odd
indices.  Hence $A\in\mathbb Z[i]^{252}$.  Antisymmetric extension
uses only permutation signs and zeros for repeated indices.
The certificate also checks the Hodge identities, the real canonical
metric, and the exact background reconstruction.

The output Sigma covectors retain their canonical normalization.
For output Phi covectors we may use $E_{ab}+E_{ba}$, $a<b$, and
$E_{aa}-E_{99}$, $a<9$; for output $H,S$ covectors use $1,i$.
These are invertible real changes of basis, so a vanishing mixed image in
these rational bases is equivalent to a vanishing canonical mixed image.

\subsection{Removing factorial sums before differentiating}

For sorted $I$ define the compressed contractions
\begin{align}
 M_{mn}(U,V)&=\sum_{|I|=4}U_{Im}V_{In},\\
 N_{ab,cd}(U,V)&=\sum_{|I|=3}U_{Iab}V_{Icd},
 \qquad a<b,\ c<d.
\end{align}
An unsorted pair is extended antisymmetrically and a repeated pair is zero.
The original normalization factors $4!$, $3!$, and the ordered-pair factors
then give exactly
\begin{align}
 I_{\lambda_2}&=\langle M(\Sigma,\bar\Sigma),M(\Sigma,\bar\Sigma)\rangle,\\
 I_{\lambda_4}&=\langle N(\Sigma,\bar\Sigma),N(\Sigma,\bar\Sigma)\rangle,\\
 I_{\lambda'_4}&=\langle N(\Sigma,\bar\Sigma),\mathcal C N(\Sigma,\bar\Sigma)\rangle,\\
 (\mathcal C N)_{ab,cd}
 &=N_{ac,bd}-N_{bc,ad}-N_{ad,bc}+N_{bd,ac}.
\end{align}
Here $\langle X,Y\rangle=\sum X_{ij}Y_{ij}$ is a bilinear dot product,
\emph{not} a Hermitian norm.  Confusing the two would erase the important
holomorphic cancellations.  The finite integer permutation matrix
$\mathcal C$ is checked to be self-adjoint for this bilinear dot product.

For either $X=M$ or $X=N$, define
\begin{align}
 \widehat X_0&=X(A,\bar A),\\
 G_j&=X(C_j,\bar A)+X(A,\bar C_j),\\
 Q_{jk}&=X(C_j,\bar C_k)+X(C_k,\bar C_j).
\end{align}
Since $B=A/4$ and $\partial_j\Sigma=C_j/2$, differentiating the
quadratic Gram functional gives the exact Gaussian-integer formula
\begin{equation}
 \boxed{32(H_\lambda)_{jk}
 =\langle G_j,\mathcal C G_k\rangle
  +\langle\widehat X_0,\mathcal C Q_{jk}\rangle.}
\end{equation}
For $\lambda_2,\lambda_4$, replace $\mathcal C$ by the identity.
The resulting imaginary parts vanish exactly.  The matrices are at most
$252\times252$; the compressed Gram arrays are $10\times10$ and
$45\times45$.

\subsection{The mixed covectors are part of the certificate}

The other nontrivial contractions with a Sigma input are
\begin{align}
 I_\beta&=\sum_{ijkl}\Phi_{ij}\Phi_{kl}
                  N_{ik,jl}(\Sigma,\bar\Sigma),\\
 I_{\chi_4}&=S\sum_{mn}\Phi_{mn}M_{mn}(\Sigma,\Sigma)+\mathrm{c.c.},\\
 I_{\eta_1}&=\frac16\sum_{a<b,c<d,n}
         N_{ab,cd}(\Sigma,\bar\Sigma)\Sigma_{abcdn}H_n
         +\mathrm{c.c.}
\end{align}
They are differentiated before setting the external $H$ background to zero.
This includes all Phi--Sigma covectors for $\beta$, all Phi--Sigma and
S--Sigma covectors for $\chi_4$, and all twenty H--Sigma covectors for
$\eta_1$.  In particular, the fact that $I_{\eta_1}=0$ at $H=0$ does
not justify dropping its mixed Hessian.
The remaining active norm terms are $\nu^2,\lambda_0,\alpha,\chi_2$.
All nineteen other invariants either have no Sigma input or retain a
factor of $H$ after the two derivatives, and therefore give zero.

\subsection{Finite exact certificates and their continuation}

The stored integer matrices $K_{50},K_{36}$, each $252\times252$,
are treated as candidate matrices, with $P_r=K_r/8$ embedded into the
Sigma block.  The following are checked using exact integer arithmetic:
\begin{equation}
 K_r^T=K_r,\qquad K_r^2=8K_r,\qquad
 \operatorname{tr}K_r=8r,\qquad K_{50}K_{36}=0.
\end{equation}
Symmetry and idempotency make the trace an exact rank certificate.
For every one of the $29$ unit invariants, its independently derived
rational $328\times252$ covector matrix $H_\alpha=N_\alpha/d_\alpha$
is checked against
\begin{equation}
 N_\alpha K_r=d_\alpha c_{\alpha,r}\,\iota_\Sigma K_r,
 \qquad r=50,36.
\end{equation}
All denominators are cleared; every residual entry is the integer zero.
At the rational background the only nonzero eigenvalues are
\begin{equation}
\begin{array}{c|rrrrrrrr}
 \text{invariant}&\nu^2&\lambda_0&\lambda_2&\lambda_4&\lambda'_4&\alpha&\beta&\chi_2\\ \hline
 c_{\alpha,50}&-\frac12&1&8&8&32&30&-30&120\\
 c_{\alpha,36}&-\frac12&1&12&12&16&30&-30&120.
\end{array}
\end{equation}
Every individual invariant Hessian block has a single radial multidegree:
the invariant degree minus the degrees of its two external legs.
Thus each zero block and scalar Sigma block proved at this nonzero radial
point extends exactly to every nonzero radial point of the same angular
orientation.  Linearity in the real action parameters then proves
\begin{align}
 m_{50}^2={}&-\frac{\nu^2}{2}
 +\left(\frac{\lambda_0}{2}+4\lambda_2+4\lambda_4+16\lambda'_4\right)\sigma^2
 +\frac65(\alpha-\beta)w^2+60\chi_2v_s^2,\\
 m_{36}^2={}&-\frac{\nu^2}{2}
 +\left(\frac{\lambda_0}{2}+6\lambda_2+6\lambda_4+8\lambda'_4\right)\sigma^2
 +\frac65(\alpha-\beta)w^2+60\chi_2v_s^2.
\end{align}
Consequently the action-to-projector assumption in the earlier
$1/\tau$ argument has been removed.  The condition that the retained hard
spectrum is positive, and the restriction to the scalar hard one-loop
kinetic kernel, are not removed.

\paragraph{Arithmetic audit.}
The implementation uses pairs of signed integer arrays for Gaussian
integers. Before each large contraction product and sum it bounds every
resulting entry by the corresponding dimension times input entry bounds.
The fixed geometry constructors separately use only $0,\pm1$, the
diagonal entries $-2,3$, and explicitly small integer factors.
The largest recorded conservative bound is $1{,}008{,}000$, safely below
$2^{62}$ and hence below signed integer overflow.  No small residual is
accepted: the certificate consists of exact zero-entry tests.  The source
action is pinned by its SHA-256 digest; a changed action requires a fresh
audit of the symbolic contraction transcription.
